\documentclass[10pt, a4paper]{article}

% --- PAQUETES PRINCIPALES ---
\usepackage[table,xcdraw,svgnames]{xcolor}
\usepackage{geometry}
\usepackage{fontspec}
\usepackage{amsmath, amsfonts, amssymb}
\usepackage{graphicx}
\usepackage[spanish]{babel}
\usepackage{titlesec}
\usepackage{fancyhdr}
\usepackage{placeins}
\usepackage{enumitem}
\usepackage{booktabs}
\usepackage{float}
\usepackage[official]{eurosym}
\usepackage{tikz}
% Añadimos la librería 'babel' para evitar conflictos con los caracteres < y >
\usetikzlibrary{shapes, arrows.meta, positioning, fit, calc, shadows, backgrounds, automata, babel}
\usepackage{tcolorbox}
\tcbuselibrary{skins, breakable}
\usepackage{fontawesome5}
\usepackage{listings}
\usepackage[colorlinks=true, linkcolor=DarkBlue, urlcolor=DarkBlue]{hyperref}

% --- CONFIGURACIÓN ---
\setmainfont{CMU Sans Serif}
\setsansfont{CMU Sans Serif}
\renewcommand{\familydefault}{\sfdefault}
\geometry{a4paper,left=2cm,right=2cm,top=2.5cm,bottom=2.5cm}

% --- DEFINICIÓN DE COLORES ---
\definecolor{MainBlue}{RGB}{0, 75, 145}
\definecolor{SoftBlue}{RGB}{235, 245, 255}
\definecolor{MainGreen}{RGB}{20, 100, 40}
\definecolor{SoftGreen}{RGB}{235, 250, 240}
\definecolor{MainOrange}{RGB}{200, 80, 0}
\definecolor{SoftOrange}{RGB}{255, 245, 230}
\definecolor{DarkBlue}{HTML}{00008B}
\definecolor{DarkRed}{HTML}{8B0000}
\definecolor{SafeZone}{RGB}{225, 250, 225}
\definecolor{PowerZone}{RGB}{250, 225, 225}
\definecolor{vhdlbg}{RGB}{248, 250, 255}
\definecolor{vhdlkw}{RGB}{0, 0, 180}
\definecolor{vhdlcmt}{RGB}{100, 100, 100}
\definecolor{vhdlstr}{RGB}{0, 128, 0}

% --- ESTILOS TIKZ ---
\tikzset{
	bloque/.style={rectangle, draw=MainBlue, thick, fill=SoftBlue,
		text width=3cm, align=center, minimum height=1.3cm,
		rounded corners=4pt, drop shadow, font=\sffamily\small},
	nucleo/.style={rectangle, draw=black!60, thick, dashed,
		fill=white!90!gray, inner sep=0.4cm, align=center, rounded corners=4pt},
	comunicacion/.style={draw=red, line width=1.5pt, <->, >=Latex},
	esquema/.style={rectangle, draw=black, thick, fill=white,
		align=center, minimum height=0.8cm},
	estado/.style={circle, draw=MainBlue, thick, fill=SoftBlue,
		minimum size=1.8cm, align=center, font=\sffamily\small\bfseries, drop shadow},
	flecha/.style={->, >=Latex, thick, color=MainBlue}
}

% --- VHDL LISTINGS ---
\lstdefinelanguage{VHDLcustom}{
	morekeywords={library, use, entity, is, port, in, out, inout, end,
		architecture, of, signal, begin, process, variable,
		if, then, else, elsif, case, when, others,
		rising_edge, falling_edge, wait, until, for,
		generic, map, all, downto, to, not, and, or, xor,
		nand, nor, type, integer, range, natural,
		std_logic, std_logic_vector, ieee},
	sensitive=false,
	morecomment=[l]{--},
	morestring=[b]"
}
\lstset{
	language=VHDLcustom,
	basicstyle=\ttfamily\small,
	keywordstyle=\bfseries\color{vhdlkw},
	commentstyle=\itshape\color{vhdlcmt},
	stringstyle=\color{vhdlstr},
	backgroundcolor=\color{vhdlbg},
	frame=single, framerule=0.5pt,
	rulecolor=\color{MainBlue!40},
	numbers=left, numberstyle=\tiny\color{gray!70},
	numbersep=8pt,
	breaklines=true, captionpos=b,
	showstringspaces=false,
	tabsize=4,
	xleftmargin=0.6cm,
	aboveskip=8pt, belowskip=6pt
}
\renewcommand{\lstlistingname}{Listado}

% --- CABECERA INTERACTIVA ---
\pagestyle{fancy}
\fancyhf{}
\setlength{\headheight}{25pt}
\fancyhead[L]{\small\textbf{\hyperlink{inicio}{Diseño Digital 2}}}
\fancyhead[R]{\small\textbf{\hyperlink{section.\thesection}{Memoria --- Fase 1}}}
\fancyfoot[C]{\thepage}
\renewcommand{\footrulewidth}{0.4pt}

% --- ESTILOS DE SECCIÓN ---
\titleformat{\section}{\Large\bfseries\color{MainBlue}}{\thesection.}{0.5em}{}
\titleformat{\subsection}{\large\bfseries\color{MainBlue}}{\thesubsection.}{0.5em}{}

% =====================================================================
\begin{document}
	\hypertarget{inicio}{}
	
	% =====================================================================
	%  PORTADA
	% =====================================================================
\begin{titlepage}
	\thispagestyle{empty}
	
	% ---- Fondo simplificado (Solo franja superior) ----
	\begin{tikzpicture}[remember picture, overlay]
		% Fondo blanco completo
		\fill[white] (current page.south west) rectangle (current page.north east);
		
		% Franja azul superior limpia (~20% de la página)
		\fill[MainBlue]
		(current page.north west) rectangle
		([yshift=-0.22\paperheight]current page.north east);
		
		% Línea de acento sutil
		\draw[MainBlue!40, line width=2pt]
		([yshift=-0.22\paperheight]current page.north west) --
		([yshift=-0.22\paperheight]current page.north east);
	\end{tikzpicture}
	
	% ---- TEXTOS DE LA ZONA AZUL ----
	\vspace*{-1cm} % Ajuste para subir el texto a la franja azul
	\begin{flushleft}
		{\fontsize{9}{12}\selectfont\color{white!80} GRADO EN INGENIERÍA ELECTRÓNICA DE COMUNICACIONES}\\[1ex]
		{\fontsize{24}{28}\selectfont\bfseries\color{white} Diseño Digital 2}\\[1ex]
		{\fontsize{11}{14}\selectfont\color{white!80} Curso académico 2025--2026}
	\end{flushleft}
	
	% ---- CONTENIDO PRINCIPAL ----
	\vspace*{0.12\paperheight}
	\begin{center}
		
		% Subtítulo / Fase
		{\Large\color{MainBlue!80}\textbf{BLOQUE II} \quad\textbullet\quad Conversores de Código}\\[1cm]
		
		% Título Principal
		{\fontsize{30}{38}\selectfont\bfseries\color{MainBlue}
			Memoria del Proyecto}\\[0.2cm]
		{\fontsize{22}{28}\selectfont\bfseries\color{MainBlue!80}
			Bloque II (Fase I)}\\[0.6cm]
		
		% Descripción técnica y versión
		{\fontsize{12}{18}\selectfont\color{black!60}
			Diseño, implementación y verificación en VHDL}\\[0.2cm]
		{\fontsize{12}{18}\selectfont\ttfamily\color{black!50}
			bcd\_to\_bin \quad$\longleftrightarrow$\quad bin\_to\_bcd}\\[0.5cm]
		
		{\fontsize{11}{14}\selectfont\bfseries\color{MainBlue}
			(Versión 1.0 --- Mayo 2026)}
		
		\vspace{1.2cm}
		
		% Tarjeta de metadatos (Simplificada sin bordes pesados)
		\renewcommand{\arraystretch}{1.5}
		\begin{tabular}{@{}r@{\hspace{1em}}l@{}}
			{\small\color{MainBlue}\faLayerGroup}\enspace\textbf{Fase:} & Bloque II (Fase I) \\
			{\small\color{MainBlue}\faCalendar}\enspace\textbf{Fecha:}  & 6 Mayo de 2026 \\
		\end{tabular}
		
		\vspace{1.2cm}
		
		% ---- SECCIÓN DE AUTORES / GRUPO ----
		\begin{tcolorbox}[
			enhanced, width=0.80\textwidth, % Ligeramente más ancho para que quepan bien los correos largos
			colback=SoftBlue, colframe=MainBlue,
			boxrule=0.8pt, arc=4pt,
			title={\centering\textbf{\faUsers\quad Miembros del Grupo}},
			coltitle=white,
			fonttitle=\large
			]
			\begin{center}
				\renewcommand{\arraystretch}{1.3}
				\begin{tabular}{@{}ll@{}}
					\textbf{Lucía Llana Carmona} & \small \texttt{lucia.llana@alumnos.upm.es} \\
					\textbf{María Moya Marcilla} & \small \texttt{maria.moya.marcilla@alumnos.upm.es} \\
					\textbf{Pedro Gil Bolaños} & \small \texttt{pedro.gilb@alumnos.upm.es} \\
					\textbf{Pedro Osma Moya} & \small \texttt{p.osma@alumnos.upm.es} \\
					\textbf{Eduardo Muñoz de Maya} & \small \texttt{eduardo.mdemaya@alumnos.upm.es} \\
				\end{tabular}
			\end{center}
		\end{tcolorbox}
		
	\end{center}
	
	\vfill
	
	% ---- PIE DE PÁGINA ----
	\begin{center}
		{\fontsize{8}{10}\selectfont\color{black!40}
			Grado en Ingeniería Electrónica de Comunicaciones \quad$\cdot$\quad Diseño Digital 2}
	\end{center}
	\vspace{0.5cm}
	
\end{titlepage}
	% =====================================================================
	%  ÍNDICE
	% =====================================================================
	\newpage
	\tableofcontents
	\newpage
	
	% =====================================================================
	\section{Introducción y Objetivos}
	% =====================================================================
	
	En esta primera fase el trabajo ha consistido en proponer el \textbf{diseño jerárquico} de la calculadora y en modelar, simular y depurar los dos conversores de código: \textbf{BCD a Binario} (\texttt{bcd\_to\_bin}) y \textbf{Binario a BCD} (\texttt{bin\_to\_bcd}).
	
	Estos dos módulos son los más importantes de cara a la interfaz del sistema. El usuario introduce y recibe los datos en BCD (teclado y display de siete segmentos), pero toda la aritmética interna ---suma, resta y multiplicación--- se hace en binario en complemento a dos. Por eso necesitamos estos conversores: actúan de puente entre los dos dominios.
	
	\vspace{0.5em}
	\begin{tcolorbox}[colback=SoftBlue, colframe=MainBlue, boxrule=0.8pt, arc=2pt,
		left=0.5cm, right=0.5cm, top=0.3cm, bottom=0.3cm]
		\small\faInfoCircle\enspace\textbf{Todos los módulos} emplean
		\texttt{ieee.std\_logic\_unsigned}. Este paquete no soporta la multiplicación
		directa de vectores de longitud arbitraria, por lo que se recurre a
		desplazamientos (shifts) y sumas, que son la representación hardware natural
		de las potencias de dos.
	\end{tcolorbox}
	
	\subsection{Arquitectura jerárquica del sistema}
	
	La calculadora se estructura de forma jerárquica con los siguientes bloques principales:
	
	\vspace{0.8em}
	\begin{center}
		\begin{tikzpicture}[node distance=0.9cm and 1.4cm]
			% Bloques principales
			\node[bloque, text width=2.6cm] (bcd2bin)
			{\texttt{bcd\_to\_bin}\\[2pt]{\tiny BCD $\to$ Binario\\combinacional}};
			\node[bloque, text width=2.6cm, right=2.0cm of bcd2bin] (alu)
			{ALU\\[2pt]{\tiny Aritmética Binaria\\suma / resta / mult.}};
			\node[bloque, text width=2.6cm, right=2.0cm of alu] (bin2bcd)
			{\texttt{bin\_to\_bcd}\\[2pt]{\tiny Binario $\to$ BCD\\secuencial}};
			
			% Bloques periféricos
			\node[bloque, text width=2.6cm, below=1.1cm of bcd2bin,
			fill=SoftGreen, draw=MainGreen] (entrada)
			{Entrada BCD\\[2pt]{\tiny Teclado / Switches}};
			\node[bloque, text width=2.6cm, below=1.1cm of bin2bcd,
			fill=SoftGreen, draw=MainGreen] (salida)
			{Salida BCD\\[2pt]{\tiny Display 7 segmentos}};
			\node[bloque, text width=2.6cm, below=1.1cm of alu,
			fill=SoftOrange, draw=MainOrange] (ctrl)
			{Control\\[2pt]{\tiny FSM Principal}};
			
			% Buses de datos
			\draw[flecha] (bcd2bin) -- node[above,font=\tiny]{10 bits} (alu);
			\draw[flecha] (alu)     -- node[above,font=\tiny]{20 bits} (bin2bcd);
			\draw[flecha] (entrada) -- (bcd2bin);
			\draw[flecha] (bin2bcd) -- (salida);
			\draw[flecha, dashed, color=MainOrange] (ctrl) -- (alu);
			
			% Marco de Fase 1
			\begin{scope}[on background layer]
				\node[nucleo, fit=(bcd2bin)(bin2bcd)(alu),
				label={[font=\small\bfseries\color{MainBlue}]above:
					\textbf{Fase 1} --- Conversores de Código}] {};
			\end{scope}
		\end{tikzpicture}
	\end{center}
	\FloatBarrier
	
	\vspace{0.4em}
	Los bloques \texttt{bcd\_to\_bin} y \texttt{bin\_to\_bcd} son, por tanto, la frontera entre el dominio decimal del usuario y el binario interno. En esta fase nos hemos centrado en diseñarlos y verificarlos con sus respectivos testbenches.
	
	% =====================================================================
	\section{Conversor BCD a Binario (\texttt{bcd\_to\_bin})}
	% =====================================================================
	
	\subsection{Especificación}
	
	\begin{tcolorbox}[colback=SoftGreen, colframe=MainGreen, boxrule=0.8pt, arc=2pt,
		left=0.5cm, right=0.5cm, top=0.3cm, bottom=0.3cm]
		\small
		\begin{tabular}{ll}
			\textbf{Tipo:}    & Lógica combinacional pura (sin reloj).\\[3pt]
			\textbf{Entrada:} & \texttt{bcd(11..0)} --- tres nibbles:
			centenas \texttt{bcd[11:8]}, decenas \texttt{bcd[7:4]},
			unidades \texttt{bcd[3:0]}.\\[3pt]
			\textbf{Salida:}  & \texttt{bin(9..0)} --- valor binario en
			$\left[0,\,999\right]$.
		\end{tabular}
	\end{tcolorbox}
	
	\vspace{0.4em}
	La conversión es simplemente:
	$$\texttt{bin} \;=\; c \times 100 \;+\; d \times 10 \;+\; u$$
	
	El máximo posible es $9 \times 100 + 9 \times 10 + 9 = 999$. Con 10 bits se pueden representar hasta $2^{10}-1 = 1023$, así que no hay riesgo de desbordamiento en ningún caso válido.
	
	\subsection{Implementación mediante desplazamientos}
	
	Como no se puede usar multiplicación directa con \texttt{std\_logic\_unsigned}, hemos descompuesto los factores 100 y 10 como suma de potencias de dos y los hemos implementado con concatenaciones de ceros (equivalente a desplazamientos a la izquierda):
	$$c \times 100 \;=\; c \times 64 + c \times 32 + c \times 4
	\;=\; (c \!\ll\! 6) + (c \!\ll\! 5) + (c \!\ll\! 2)$$
	$$d \times 10 \;=\; d \times 8 + d \times 2
	\;=\; (d \!\ll\! 3) + (d \!\ll\! 1)$$
	
	\begin{lstlisting}[caption={Núcleo de \texttt{bcd\_to\_bin.vhd} --- multiplicación por desplazamientos}]
		-- Extender los tres digitos BCD a 10 bits cada uno
		c <= "000000" & bcd(11 downto 8);   -- centenas
		d <= "000000" & bcd( 7 downto 4);   -- decenas
		u <= "000000" & bcd( 3 downto 0);   -- unidades
		
		-- c * 100 = c*64 + c*32 + c*4  (64+32+4 = 100)
		c100 <= (c(3 downto 0) & "000000")   -- c * 64  (shift 6)
		+ (c(4 downto 0) & "00000")    -- c * 32  (shift 5)
		+ (c(7 downto 0) & "00");      -- c * 4   (shift 2)
		
		-- d * 10 = d*8 + d*2  (8+2 = 10)
		d10  <= (d(6 downto 0) & "000")      -- d * 8   (shift 3)
		+ (d(8 downto 0) & '0');       -- d * 2   (shift 1)
		
		bin <= c100 + d10 + u;
	\end{lstlisting}
	
	\begin{tcolorbox}[colback=SoftOrange, colframe=MainOrange, boxrule=0.8pt, arc=2pt,
		left=0.5cm, right=0.5cm, top=0.3cm, bottom=0.3cm]
		\small\faExclamationTriangle\enspace\textbf{Anchos de bits:} Cada
		concatenación produce exactamente 10 bits, igual que la señal destino.
		Por ejemplo, \texttt{c(3 downto 0) \& "000000"} produce
		$4 + 6 = 10$ bits. Con \texttt{std\_logic\_unsigned} las sumas operan
		en ese ancho fijo, descartando cualquier carry sin riesgo de
		desbordamiento silencioso.
	\end{tcolorbox}
	
	\subsection{Verificación}
	
	El testbench \texttt{tb\_bcd\_to\_bin} simplemente va cambiando la entrada cada 20\,ns sin necesidad de reloj, ya que el módulo es puramente combinacional. Los casos cubren los extremos y algunos valores intermedios:
	
	\begin{table}[H]
		\centering
		\caption{Casos de prueba del testbench \texttt{tb\_bcd\_to\_bin}}
		\begin{tabular}{ccc}
			\toprule
			\textbf{Entrada BCD (hex)} & \textbf{Salida binaria (decimal)} & \textbf{Propósito} \\
			\midrule
			\texttt{0x000} & 0   & Cero absoluto \\
			\texttt{0x010} & 10  & Solo decenas  \\
			\texttt{0x100} & 100 & Solo centenas \\
			\texttt{0x123} & 123 & Valor mixto representativo \\
			\texttt{0x999} & 999 & Máximo absoluto \\
			\bottomrule
		\end{tabular}
		\label{tab:tb-bcd2bin}
	\end{table}
	
	\FloatBarrier
	
	% =====================================================================
	\section{Conversor Binario a BCD (\texttt{bin\_to\_bcd})}
	% =====================================================================
	
	\subsection{Especificación}
	
	\begin{tcolorbox}[colback=SoftGreen, colframe=MainGreen, boxrule=0.8pt, arc=2pt,
		left=0.5cm, right=0.5cm, top=0.3cm, bottom=0.3cm]
		\small
		\begin{tabular}{ll}
			\textbf{Tipo:}      & Lógica secuencial --- FSM de tres estados.\\[3pt]
			\textbf{Entradas:}  & \texttt{clk} (50\,MHz), \texttt{nRst} (reset activo bajo),
			\texttt{start}, \texttt{bin(N\_BITS-1..0)}.\\[3pt]
			\textbf{Salidas:}   & \texttt{done} (pulso de un ciclo),
			\texttt{bcd(23..0)} --- seis dígitos BCD.\\[3pt]
			\textbf{Genérico:}  & \texttt{N\_BITS = 20} $\;\Rightarrow\;$
			rango $\left[0,\, 2^{20}-1\right] = \left[0,\,1\,048\,575\right]$.
		\end{tabular}
	\end{tcolorbox}
	
	\vspace{0.4em}
	La calculadora tiene que poder mostrar resultados de hasta $999 \times 999 = 998\,001$, así que el conversor trabaja con 20 bits de entrada y produce 6 dígitos BCD (24 bits). Hacer esto con lógica combinacional pura sería inviable, así que hemos optado por una \textbf{arquitectura secuencial iterativa}.
	
	\subsection{Algoritmo: Suma de pesos (Horner en BCD)}

	Hemos implementado el algoritmo de suma de pesos, que construye el resultado BCD
	procesando el número binario bit a bit de más significativo a menos, en exactamente
	$N_{\text{BITS}} = 20$ ciclos de reloj. La idea es aplicar el \textit{método de Horner}
	pero realizando la aritmética directamente en BCD:

	$$\text{acc} \leftarrow 2 \cdot \text{acc} + b_i \quad \text{(aritmética BCD)}, \qquad i = N-1, N-2, \ldots, 0$$

	Cada ciclo hace lo siguiente:

	\begin{enumerate}[leftmargin=1.8cm, label=\textbf{Paso \arabic*.}]
		\item \textbf{Doblar en BCD}: para cada dígito BCD del acumulador se calcula
		$\mathit{tmp} = 2 \cdot d_k + c_{\text{entrada}}$ (5~bits, rango 0--19). Si
		$\mathit{tmp} \geq 10$, se resta 10 al resultado y se genera $c_{\text{salida}}=1$
		que se propaga al dígito superior.

		\item \textbf{Sumar el bit entrante}: el bit MSB de \texttt{bin\_reg} actúa
		como $c_{\text{entrada}}$ del dígito de unidades (vale 0 o 1, por lo que nunca
		causa desbordamiento por sí solo).
	\end{enumerate}

	\begin{tcolorbox}[colback=SoftBlue, colframe=MainBlue, boxrule=0.8pt, arc=2pt,
		left=0.4cm, right=0.4cm, top=0.2cm, bottom=0.2cm]
		\small\faLightbulb\enspace\textbf{Ejemplo (8 bits, $10011101_2 = 157_{10}$):}
		Peso$_{128}$=1,\ Peso$_{64}$=$2{\times}1{+}0=2$,\ Peso$_{32}$=$2{\times}2{+}0=4$,\
		Peso$_{16}$=$2{\times}4{+}1=9$,\ Peso$_8$=$2{\times}9{+}1=19$,\
		Peso$_4$=$2{\times}19{+}1=39$,\ Peso$_2$=$2{\times}39{+}0=78$,\
		Peso$_1$=$2{\times}78{+}1=157$ (\textit{todo en BCD})
	\end{tcolorbox}

	\begin{lstlisting}[caption={Núcleo del algoritmo de suma de pesos en \texttt{bin\_to\_bcd.vhd}}]
		-- tmp = 2*digito + carry_in  (5 bits, rango 0..19)
		-- Si tmp >= 10: resultado = tmp(3:0) + 6  (resta 10 = suma 6 en nibble)
		--               carry_out = 1
		c := bin_reg(N_BITS-1);   -- bit entrante -> carry inicial del digito 0
		tmp := ('0' & acc_bcd(3 downto 0)) + ('0' & acc_bcd(3 downto 0)) + ("0000" & c);
		if tmp >= "01010" then n0 := tmp(3 downto 0) + "0110"; c := '1';
		else n0 := tmp(3 downto 0); c := '0'; end if;
		-- (idem para n1..n5: decenas, centenas, miles, decenas de miles, centenas de miles)

		acc_bcd   <= n5 & n4 & n3 & n2 & n1 & n0;
		bin_reg   <= bin_reg(N_BITS-2 downto 0) & '0';  -- avanzar registro binario
		bits_left <= bits_left - 1;
	\end{lstlisting}
	
	\subsection{Máquina de Estados (FSM)}
	
	El control secuencial se implementa mediante una FSM de tres estados:
	
	\vspace{0.8em}
	\begin{center}
		\begin{tikzpicture}[
			>=Latex,
			every node/.style={font=\sffamily\small}
			]
			% ---- Nodos en triángulo ----
			% IDLE  abajo-izquierda, CALC  arriba-centro, FINISH  abajo-derecha
			\node[estado] (idle)   at (0,0)    {IDLE};
			\node[estado] (calc)   at (4,3)    {CALC};
			\node[estado] (finish) at (8,0)    {FINISH};
			
			% ---- IDLE self-loop (lado izquierdo del nodo) ----
			\draw[flecha] (idle) to[out=210, in=150, looseness=6]
			node[left, font=\tiny\sffamily, align=right, xshift=-2pt]
			{\texttt{start='0'}}
			(idle);
			
			% ---- IDLE -> CALC ----
			\draw[flecha] (idle) to[bend left=12]
			node[above left, font=\tiny\sffamily, align=center, xshift=2pt]
			{\texttt{start='1'}\\[-1pt]carga \texttt{bin\_reg}}
			(calc);
			
			% ---- CALC self-loop (arriba del nodo) ----
			% Corrección: Sacamos el símbolo '>' del \texttt para evitar conflicto activo de babel
			\draw[flecha] (calc) to[out=60, in=120, looseness=6]
			node[above, font=\tiny\sffamily, align=center]
			{	\texttt{bits\_left} $>$ 0\\[-1pt]Horner BCD}
			(calc);
			
			% ---- CALC -> FINISH ----
			\draw[flecha] (calc) to[bend left=12]
			node[above right, font=\tiny\sffamily, align=center, xshift=-2pt]
			{\texttt{bits\_left = 0}}
			(finish);
			
			% ---- FINISH -> IDLE (arco inferior) ----
			\draw[flecha] (finish) to[bend left=25]
			node[below, font=\tiny\sffamily, align=center]
			{\texttt{done} $\leftarrow$ \texttt{'1'}}
			(idle);
			
		\end{tikzpicture}
	\end{center}
	\FloatBarrier
	
	\begin{table}[H]
		\centering
		\caption{Estados de la FSM del conversor \texttt{bin\_to\_bcd}}
		\begin{tabular}{lll}
			\toprule
			\textbf{Estado} & \textbf{Descripción} & \textbf{Acción sobre señales} \\
			\midrule
			\texttt{IDLE}
			& Espera activa de \texttt{start='1'}.
			Carga \texttt{bin\_reg} $\leftarrow$ \texttt{bin}.
			& \texttt{done} $\leftarrow$ \texttt{'0'} \\[4pt]
			\texttt{CALC}
			& Ejecuta un paso Horner BCD por ciclo
			durante \texttt{N\_BITS} ciclos.
			& \texttt{acc\_bcd}, \texttt{bin\_reg} actualizados \\[4pt]
			\texttt{FINISH}
			& Pulso de fin; retorna inmediatamente a \texttt{IDLE}.
			& \texttt{done} $\leftarrow$ \texttt{'1'} \\
			\bottomrule
		\end{tabular}
		\label{tab:fsm}
	\end{table}
	
	\FloatBarrier
	
	\subsection{Verificación}
	
	El testbench \texttt{tb\_bin\_to\_bcd} genera un reloj de 50\,MHz (periodo 20\,ns). Nos encontramos con problemas de condiciones de carrera al detectar \texttt{done='1'}, así que lo resolvimos con un procedimiento auxiliar (\texttt{run\_test}) que alinea todo con flancos de reloj. La simulación termina automáticamente mediante \texttt{assert false severity failure}, por lo que basta con ejecutar \texttt{run -all}.
	
	\begin{lstlisting}[caption={Procedimiento auxiliar \texttt{run\_test} del testbench secuencial}]
		procedure run_test(val : std_logic_vector(19 downto 0)) is
		begin
		wait until clk'event and clk = '1';  -- alinear con el reloj antes de asignar
		bin   <= val;
		start <= '1';
		wait until clk'event and clk = '1';  -- IDLE detecta start; done pasa a '0'
		start <= '0';
		wait until done = '1';               -- esperar fin de la conversion (20 ciclos)
		wait for clk_period * 2;             -- margen para visualizar en waveform
		end procedure;
	\end{lstlisting}
	
	\begin{table}[H]
		\centering
		\caption{Casos de prueba del testbench \texttt{tb\_bin\_to\_bcd}}
		\begin{tabular}{ccl}
			\toprule
			\textbf{Binario (decimal)} & \textbf{BCD esperado (hex)} & \textbf{Propósito} \\
			\midrule
			0         & \texttt{0x000000} & Cero absoluto \\
			1         & \texttt{0x000001} & Unidad mínima \\
			9         & \texttt{0x000009} & Máximo de un dígito \\
			10        & \texttt{0x000010} & Decena exacta \\
			100       & \texttt{0x000100} & Centena exacta \\
			255       & \texttt{0x000255} & Potencia de 2 menos 1 \\
			999       & \texttt{0x000999} & Máximo operando de entrada \\
			1\,000    & \texttt{0x001000} & Paso de millar \\
			9\,999    & \texttt{0x009999} & Máximo de cuatro dígitos \\
			999\,999  & \texttt{0x999999} & Máximo BCD representable (6 dígitos) \\
			\bottomrule
		\end{tabular}
		\label{tab:tb-bin2bcd}
	\end{table}
	
	\noindent
	\textit{Nota: el valor máximo operacional de la calculadora es $999 \times 999 = 998\,001$.
		El último caso de prueba utiliza $999\,999$ porque es el máximo que los 24 bits BCD de
		salida pueden representar, ofreciendo así una verificación de frontera más exigente.}
	
	\FloatBarrier
	
	% =====================================================================
	\section{Conclusiones}
	% =====================================================================
	
	En esta primera fase hemos diseñado, implementado y verificado los dos conversores de código que forman la base del proyecto:
	
	\begin{itemize}[leftmargin=1.6cm, itemsep=4pt]
		\item \textbf{\texttt{bcd\_to\_bin}}: completamente combinacional. Hemos resuelto la multiplicación por 100 y por 10 mediante desplazamientos y sumas, sin usar multiplicadores, lo que lo hace compatible con \texttt{std\_logic\_unsigned}. La latencia es nula (solo depende de la propagación de puertas).
		
		\item \textbf{\texttt{bin\_to\_bcd}}: secuencial con FSM de tres estados (\texttt{IDLE} / \texttt{CALC} / \texttt{FINISH}). El algoritmo de suma de pesos (Horner en BCD) tarda exactamente 20 ciclos de reloj, lo que es más que suficiente para los valores que puede generar la calculadora ($999 \times 999 = 998\,001$ como máximo).
	\end{itemize}
	
	Todos los casos de prueba han dado el resultado esperado, incluyendo el cero, los valores de frontera y algunos intermedios. Los módulos quedan listos para integrarse en las siguientes fases.
	
	\vspace{1em}
	\begin{tcolorbox}[colback=SoftBlue, colframe=MainBlue, boxrule=0.8pt, arc=2pt,
		left=0.5cm, right=0.5cm, top=0.3cm, bottom=0.3cm]
		\small\faArrowRight\enspace\textbf{Próxima fase:}
		En la Fase 2 habrá que diseñar la ALU: suma, resta y multiplicación de operandos
		con signo en complemento a dos (10 bits), con detección de desbordamiento
		mediante la señal \texttt{overflow}.
	\end{tcolorbox}

	% =====================================================================
	\newpage
	\appendix
	\section{Investigación y conocimientos adquiridos durante el desarrollo}
	% =====================================================================

	Durante la realización de esta fase surgieron diversas dificultades técnicas cuya resolución requirió investigación adicional fuera del contenido impartido en clase. Se documentan a continuación los aspectos más relevantes.

	\subsection{Algoritmo de suma de pesos (Horner en BCD)}

	Para la conversión binario a BCD se buscaba un método que no requiriese divisiones
	ni operaciones módulo, ya que no son sintetizables directamente con
	\texttt{std\_logic\_unsigned} y resultan costosas en área hardware.

	El algoritmo elegido se basa en la \textbf{representación polinómica} de un número
	binario, evaluada mediante el método de Horner:
	$$\text{bin} = b_{N-1} \cdot 2^{N-1} + \cdots + b_1 \cdot 2 + b_0
	= (\cdots((b_{N-1}) \cdot 2 + b_{N-2}) \cdot 2 + \cdots) \cdot 2 + b_0$$

	La clave es que esta evaluación se realiza \textit{íntegramente en aritmética BCD}:
	en cada paso se dobla el acumulador BCD y se suma el siguiente bit. ``Doblar en BCD''
	significa multiplicar cada dígito por 2, y si el resultado supera 9 se resta 10 y se
	propaga un carry al dígito superior.

	\begin{tcolorbox}[colback=SoftOrange, colframe=MainOrange, boxrule=0.8pt, arc=2pt,
		left=0.5cm, right=0.5cm, top=0.3cm, bottom=0.3cm]
		\small\faLightbulb\enspace\textbf{Intuición del doblado BCD:}
		Dado un dígito $d \in [0,9]$, al doblar $d' = 2d + c_{\text{in}}$ (rango 0--19).
		Si $d' \geq 10$: el resultado decimal es $d'-10$ con $c_{\text{out}}=1$ hacia el
		dígito superior. En hardware se implementa como: si $d' \geq 10$, suma~6 al nibble
		(porque $16 - 10 = 6$, lo que equivale a restar 10 aprovechando el desbordamiento
		natural de 4~bits) y activa el carry.
	\end{tcolorbox}

	La comprensión del algoritmo requirió trazar paso a paso el ejemplo de 8~bits con
	el número $10011101_2 = 157_{10}$, verificando que en cada iteración el acumulador
	BCD es exactamente el peso del bit ya procesado. Una vez comprendido, la
	implementación en VHDL resultó directa.

	\subsection{Multiplicación sin operador de multiplicación}

	Para el conversor BCD a binario fue necesario calcular $c \times 100$ y $d \times 10$ sin disponer de un operador de multiplicación genérico. La solución se obtuvo descomponiendo cada factor como suma de potencias de dos:

	$$100 = 2^6 + 2^5 + 2^2, \qquad 10 = 2^3 + 2^1$$

	Multiplicar por una potencia de dos equivale a un desplazamiento a la izquierda, que en VHDL se implementa mediante concatenación de ceros, sin coste alguno en lógica combinacional. Este enfoque no era conocido de antemano y se dedujo a partir del análisis de la representación binaria de los factores.

	
\end{document}